% Generated draft: requires mathematical and source review.
\documentclass[11pt]{article}
\usepackage[margin=25mm]{geometry}
\usepackage{fontspec}
\setmainfont{Latin Modern Roman}
\usepackage{amsmath,amssymb,mathtools,unicode-math}
\setmathfont{Latin Modern Math}
\usepackage{enumitem,needspace,xurl,hyperref}
\setlength{\parindent}{0pt}
\setlength{\parskip}{5pt}
\setlength{\emergencystretch}{3em}
\newcommand{\N}{\mathbb{N}}
\newcommand{\Z}{\mathbb{Z}}
\newcommand{\Q}{\mathbb{Q}}
\newcommand{\R}{\mathbb{R}}
\newcommand{\C}{\mathbb{C}}
\newcommand{\F}{\mathbb{F}}
\newcommand{\A}{\mathbb{A}}
\newcommand{\PP}{\mathbb{P}}
\newcommand{\E}{\mathbb{E}}
\newcommand{\eps}{\varepsilon}
\newcommand{\dd}{\,\mathrm{d}}
\begin{document}
\section*{UnsolvedMath ID: 30006265}
\textbf{Source problem code:} catalogue entry
\subsection{Definitions and mathematical statement}
The supplied record calls \(t(\lambda)\) the \emph{Steinberg quotient}, indexed by dominant weights \(\lambda\), but supplies neither its defining formula nor the ambient representation-theoretic setting. In particular, it does not specify the group \(G\), the base field or its characteristic, the weight lattice, the choice determining dominance, or any restrictions on these data. Consequently, a fully specified mathematical statement cannot be recovered from this record alone.

Write \(X^+\) provisionally for the unspecified set of dominant weights intended in the source. The stated task is: for every \(\lambda\in X^+\), identify a natural representation-theoretic or geometric object whose character, in an appropriate specified sense, equals \(t(\lambda)\). In particular, determine whether
\[
\forall\lambda\in X^+\;\exists\text{ a finite-dimensional }G\text{-module }V_\lambda
\quad\text{such that}\quad
\operatorname{ch}(V_\lambda)=t(\lambda).
\]
Here \(G\), its module category, and the character map remain to be specified from the original source; the display records the intended quantifiers and does not supply the missing hypotheses.

For clarity, if the intended meaning is the ordinary formal character of a representation over a field \(k\), computed with respect to a torus \(T\subseteq G\), its usual expression is
\[
\operatorname{ch}_T(V)=\sum_{\mu\in X^*(T)}\dim_k(V_\mu)e^\mu,
\]
where \(X^*(T)\) is the character lattice, \(V_\mu\) is the weight space of weight \(\mu\), and \(e^\mu\) is the corresponding basis element of \(\Z[X^*(T)]\). This explains a possible character convention conditionally; the supplied record does not establish that this is its precise setting or specify how \(t(\lambda)\) is placed in this ring.

The supplied dated assessment reports a geometric and categorical realization through Smith--Treumann localization and parity-sheaf stalk data, while stating that identification with the ordinary character of a finite-dimensional \(G\)-module is not established. No precise realization theorem or hypotheses are supplied. This qualification is retained as supplied curation, not asserted here as a verified mathematical result.
\subsection{Short English statement}
For every dominant weight \(\lambda\), seek a natural object with character equal to the Steinberg quotient \(t(\lambda)\), and determine whether such a character comes from a finite-dimensional \(G\)-module. The record omits the definition of \(t(\lambda)\) and the ambient setting. Its supplied dated assessment reports a geometric realization but leaves the ordinary module-character interpretation unresolved; that assessment requires independent verification.
\subsection{Sources}
\begin{itemize}
\item Supplied local UnsolvedMath registry record, permanent ID \(30006265\), source problem code \texttt{catalogue entry}; project registry \texttt{data/unsolvedmath/problems.json}. Title: Representation-Theoretic Meaning of Steinberg Quotient Characters. The record labels the status \texttt{partially\_solved} and the statement status \texttt{corrected\_verified}; these are supplied registry labels, not independent verification.
\item Source attribution: Pramod N. Achar, joint work with Paul Sobaje, ``Smith--Treumann localization and the Steinberg module.'' The supplied original question is ``What is \(t(\lambda)\) the character of?'' The registry gives proposed year 2025 and does not separately identify a proposer.
\item \emph{Enveloping Algebras and Geometric Representation Theory} (2025), Oberwolfach Reports, pp.~1271--1332, DOI 10.4171/owr/2025/25: \url{https://doi.org/10.4171/owr/2025/25}. The supplied background describes Oberwolfach Reports as documenting mathematical workshops and their open-problem discussions.
\item Supplied literature source, ``OWR 2025 report'': \url{https://ems.press/content/serial-article-files/51851}.
\item Supplied literature source, ``Achar--Sobaje, Steinberg quotients and localization'': \url{https://arxiv.org/abs/2606.21697}.
\item Supplied literature source, ``Sobaje, Steinberg quotients and Weyl characters'': \url{https://arxiv.org/abs/2301.11465}.
\item Literature assessment date: \texttt{2026-08-22}. The supplied assessment states that Smith--Treumann localization gives a geometric and categorical realization through parity-sheaf stalk data, but that identification with the ordinary character of a finite-dimensional \(G\)-module is not established. This is supplied dated curation derived from the OWR record, not independently verified here; time-sensitive status and cited claims require independent verification.
\end{itemize}
\end{document}
